\begin{textsample}{POS Dim 4 – 2016\_c5 – Score 2.00 – 2016\_c5/enade\_2016\_tec\_agro\_1c1f311707.txt}  \label{ex:f4_pos_109}
CAPÍTULO 5 – DISTRIBUIÇÃO DOS CONCEITOS 5.1 - PANORAMA NACIONAL DA DISTRIBUIÇÃO DOS CONCEITOS O Conceito Enade é um indicador de qualidade que avalia o desempenho dos estudantes a partir dos resultados do Enade.
Os valores possíveis vão de 1 ( pior situação ) a 5 ( melhor situação ).
As Unidades de \textbf{Observação} com apenas um ou sem nenhum concluinte participante não obtêm o Conceito Enade, ficando Sem Conceito ( SC ).
A Tabela 5.1 apresenta a quantidade e distribuição de cursos de Tecnologia em Agronegócio participantes do Enade/2016, por faixa de conceito e Grande Região.
Enfatiza -se, mais uma vez, que a diferença entre os cursos tabulados neste capítulo e no capítulo 2 corresponde aos cursos sem conceito, ou seja àqueles, em princípio, sem alunos concluintes que houvessem participado da prova ou àqueles que tivessem tão somente um aluno concluinte19.
Observando -se os dados da Tabela 5.1, nota -se que, dos 63 cursos participantes, 25 ( 39,7 \% ) classificaram -se com conceito 3, o valor modal.
Esse foi também o conceito modal em todas as regiões.
O conceito 4 foi o segundo mais frequente em nível nacional ( 28,6 \%, correspondendo a 18 cursos ) e o conceito 2, o terceiro ( 19,0 \%, correspondendo a 12 cursos ).
Houve, ainda, cinco cursos ( 7,9 \% ) que receberam conceito 5 e dois cursos ( 3,2 \% ) que receberam conceito 1.
Dos 68 cursos de Tecnologia em Agronegócio, um ( 1,6 \% ) ficou sem conceito ( SC ). 19 Estes cursos com somente um concluinte não foram considerados no capítulo 2.
A região Norte participou com quatro cursos ou 6,3 \% do total nacional.
Estes cursos ficaram distribuídos entre os conceitos 1, 2, 3 e 4 ( um curso, cada ).
Nenhum curso ficou sem conceito ( SC ) ou com conceito 5 nessa região.
A região Nordeste participou com cinco cursos ou 7,9 \% do total nacional.
Nota -se que tais cursos apresentaram distribuição bimodal, uma vez que as maiores frequências se concentraram em dois conceitos, os conceitos 2 e 3, atribuídos, cada qual, a dois cursos ( 40,0 \% cada ).
Ainda nessa região, um curso ( 20,0 \% ) foi avaliado com o conceito 4.
Nenhum curso ficou sem conceito ( SC ) ou com conceito 1 ou 5.
Dos 22 cursos participantes da região Sudeste ( 34,9 \% do total nacional ), o conceito 3 foi o conceito modal, como já anteriormente assinalado, tendo sido obtido por 11 cursos ( 50,0 \% ).
O conceito 2 foi atribuído a três cursos ( 13,6 \% ) e o conceito 4, a oito cursos ( 36,4 \% ).
Nessa região, nenhum curso ficou sem conceito ( SC ) ou recebeu conceito 1 ou 5.
Os 17 cursos da região Sul corresponderam a 27,0 \% do total nacional.
Houve predominância dos conceitos 3 e 4, atribuídos a 29,4 \% dos cursos da região cada ( conceitos modais ), ou seja, atribuído a cinco cursos cada na região Sul.
Os conceitos 1 e 2 foram atribuídos a dois cursos ( um curso, cada ) e o conceito 5, a quatro cursos ( 23,5 \% ).
Um curso ficou sem conceito ( SC ) na região Sul.
Já dos 15 cursos participantes na região Centro-Oeste ( 23,8 \% do total nacional ), seis receberam conceito 3, o conceito modal ( 40,0 \% ), como já fora destacado.
A cinco cursos foi atribuído conceito 2, a três cursos, conceito 4 e a um único curso, conceito 5.
Nessa região, nenhum curso ficou sem conceito ( SC ) ou recebeu conceito 1.
O Gráfico 5.1 apresenta a distribuição do Conceito Enade segundo Grande Região : poligonais mais à esquerda representam Grandes Regiões com pior distribuição, e poligonais mais à direita, Grandes Regiões com melhores conceitos.
A região Norte ( linha azul ) apresenta os cursos com conceitos mais baixos, ao passo que as regiões Sudeste ( linha vermelha ) e Sul ( linha roxa ) apresentam as distribuições com valores maiores.
Cursos em IES situadas nas regiões Nordeste ( linha verde ) e Centro-Oeste ( linha preta ) apresentam uma situação intermediária. 5.2 - CONCEITOS POR CATEGORIA ADMINISTRATIVA, POR MODALIDADE DE ENSINO E GRANDE REGIÃO A Tabela 5.2 apresenta a distribuição dos cursos participantes do Enade/2016 de Tecnologia em Agronegócio, por Categoria Administrativa e Modalidade de Ensino, de acordo com os conceitos por eles alcançados, segundo as Grandes Regiões brasileiras.
Dos 63 cursos participantes, 34 ( 54,0 \% ) eram ministrados em Instituições Privadas e 29 ( 46,0 \% ), em Instituições Públicas.
De acordo com as informações da Tabela 5.2, em termos nacionais, entre Instituições Públicas e Privadas, observa -se que, dos cinco cursos avaliados com conceito 5, apenas um era oferecido em IES Públicas e quatro, em IES Privadas.
Dos 29 cursos participantes de IES Públicas, o conceito 3 foi o valor modal, atribuído a 12 cursos ( 41,4 \% ).
Entre os demais cursos de Instituições Públicas participantes, dois cursos ( 6,9 \% ) foram avaliados com conceito 1, quatro cursos ( 13,8 \% da categoria ) obtiveram conceito 2, nove cursos ( 31,0 \% ) receberam conceito 4 e, como já comentado, um curso ( 3,4 \% ) receberam conceito 5.
Nesta categoria, um curso ( 3,4 \% ) ficou sem conceito ( SC ).
Na rede privada, o conceito modal, também, foi 3, atribuído a 13 cursos, o correspondente a 38,2 \% dos 34 cursos da categoria.
Entre os demais cursos participantes, oito ( 23,0 \% ) receberam conceito 2.
O conceito 4 foi atribuído a nove cursos ( 26,5 \% ) e o conceito 5, a quatro cursos ( 11,8 \% ).
Nesta Categoria Administrativa.
Nenhum curso ficou sem conceito ( SC ) ou com conceito 1.
Os dados referentes à distribuição dos cursos segundo a Modalidade de Ensino revelam que apenas cinco dos 63 cursos de Tecnologia em Agronegócio participantes do Enade/2016 foram oferecidos na modalidade Educação a Distância, equivalente a 7,9 \% dos cursos.
Quatro desses cursos foram avaliados com conceitos 4 e 5 ( dois cursos cada, conceitos modais ) e um curso recebeu conceito 2.
Os cursos oferecidos na modalidade a distância estavam concentrados em duas regiões : Sudeste com três cursos e Sul com dois cursos.
Os dois cursos da região Sul receberam conceito 5, enquanto os da região Sudeste receberam conceitos 2 ( um curso ) e 4 ( dois cursos ).
Na análise por região, observa -se que, na região Norte, as Instituições Privadas participaram com dois cursos ( 50,0 \% do total regional ), dos quais um ficou com conceito 2 e um com conceito 4.
As Instituições Públicas, também, participaram com dois cursos na região Norte ( 50,0 \% do total regional ), dos quais um obteve conceito 1 e o outro com conceito 3.
Na região Nordeste, a rede privada apresentou dois cursos participantes, o equivalente a 40,0 \% do total da região, a menor das proporções desta categoria entre as regiões.
Estes dois cursos foram avaliados com conceito 3.
As Instituições Públicas da região Nordeste participaram com três cursos ( 60 \% do total da região ), a maior das proporções desta categoria nas regiões.
Desses, dois obtiveram conceito 2 e um recebeu conceito 4.
Na região Sudeste, a proporção de cursos da rede privada foi de 45,5 \%, correspondendo a dez dos 22 cursos participantes.
Nessa categoria e região, o conceito modal foi 4, atribuído a quatro cursos.
Seis cursos receberam conceitos 2 e 3 ( três cursos, cada ).
Nenhum curso ficou sem conceito ( SC ) ou com os demais conceitos.
Entre os 12 cursos oferecidos por Instituições Públicas na região Sudeste, o conceito modal foi 3, atribuído a oito cursos.
Os demais cursos receberam o conceito 4 ( quatro cursos ).
Nenhum curso ficou sem conceito ( SC ) ou recebeu os demais conceitos.
As Instituições Privadas concentraram 11 dos 17 cursos participantes da região Sul, 64,7 \% do total regional, mais elevada do que a registrada nas demais regiões brasileiras.
Desses, oito obtiveram os conceitos modais 3 e 5.
Nesta combinação de Categoria Administrativa e Grande Região, os conceitos 2 e 4 foram atribuídos a um e a dois cursos, respectivamente.
Nenhum recebeu conceito 1 ou ficou sem conceito ( SC ).
As Instituições Públicas na região Sul participaram com seis cursos, 35,3 \% dos cursos da região.
O conceito modal foi 4, atribuído a três cursos.
Aos demais cursos foram atribuídos os conceitos 1 e 3 ( um curso, cada ) e um curso ficou sem conceito ( SC ).
Nenhum curso da região Sul oferecido por IES Pública recebeu os conceitos 2 e 5.
Na região Centro-Oeste, nove dos 15 cursos participantes eram de Instituições Privadas ( 60,0 \% em termos regionais ).
Desses, quatro receberam conceito 3, o conceito modal.
Entre os demais, três receberam conceito 2 e dois receberam conceito 4.
Nenhum curso ficou sem conceito ( SC ) ou com os demais conceitos.
Quanto aos seis cursos oferecidos por Instituições Públicas na região Centro-Oeste ( 40,0 \% do total regional ), nota -se que tais cursos apresentaram distribuição bimodal, uma vez que as maiores frequências se concentraram em dois conceitos, os conceitos 2 e 3, atribuídos, cada qual, a dois cursos.
Os dois cursos restantes receberam conceitos 4 e 5 ( um curso, cada ).
Nenhum curso ficou sem conceito ( SC ) ou com conceito 1.
Como já comentado cinco cursos foram oferecidos na modalidade a distância.
Três deles estavam situados na região Sudeste, avaliados com conceito 2 ( um curso ) e 4 ( dois cursos ).
Os outros dois estavam situados na região Centro-Oeste, estes avaliados com conceito 5.
O Gráfico 5.2 apresenta a distribuição do Conceito Enade segundo a Categoria Administrativa da IES.
Os cursos em IES Públicas ( linha azul ) e IES Privadas ( linha verde ) apresentam poligonais semelhantes e, consequentemente, uma distribuição de Conceitos Enade com valores parecidos. 5.3 - CONCEITOS POR ORGANIZAÇÃO ACADÊMICA E GRANDE REGIÃO Na Tabela 5.3, encontra -se a distribuição dos conceitos atribuídos aos cursos participantes do Enade/2016 na Área de Tecnologia em Agronegócio, por Organização Acadêmica, segundo as Grandes Regiões brasileiras.
Dos 63 cursos de Tecnologia em Agronegócio participantes, 27 eram oferecidos em Universidades, cinco Universitários e 31 em Faculdades.
Esta distribuição corresponde a, respectivamente, 42,9 \%, 7,9 \% e 49,2 \% dos cursos.
De acordo com os dados apresentados, dos cinco cursos avaliados com conceito 5, dois eram vinculados a Universidades.
Os cursos vinculados a esse tipo de Organização Acadêmica tiveram os conceitos 3 e 4 como conceitos modais, atribuídos a nove cursos cada ( 33,3 \%, cada ).
Os demais cursos vinculados a Universidades receberam os conceitos 1 ( um curso ), 2 ( cinco cursos ) e 5 ( dois cursos ).
Ficou sem conceito ( SC ), um curso.
Entre os cursos em Centros Universitários, o conceito modal foi 2, atribuído a três cursos ( 60,0 \% ).
Nenhum curso vinculado a esse tipo de Organização Acadêmica ficou sem conceito ( SC ) ou com conceito 1 ou 4.
Os outros cursos receberam os conceitos 3 e 5 ( um curso, cada ).
Dos 31 cursos mantidos por Faculdades, 15 ( 48,4 \% ) receberam conceito 3, o conceito modal.
O segundo conceito mais frequentemente atribuído aos cursos oferecidos por essa categoria de IES foi o conceito 4, recebido por nove cursos.
Na sequência, quatro cursos obtiveram conceito 2 e dois cursos, conceito 5.
O conceito 1 foi atribuído a um curso, e nenhum curso ficou sem conceito ( SC ).
Considerando -se separadamente as regiões brasileiras, verifica -se que, na região Norte, as Universidades concentraram dois dos quatro cursos participantes.
Um foi avaliado com conceito 1 e o outro com conceito 3.
Os Centros Universitários da região Norte foram representados por um único curso, ao qual foi atribuído o conceito 2.
As Faculdades participaram, também, com somente um curso na região Norte, e este recebeu o conceito 4.
Na região Nordeste, as Universidades participaram com três dos cinco cursos da Área de Tecnologia em Agronegócio oferecidos na região, dos quais dois ficaram com conceito 2 e um com conceito 4.
Os Centros Universitários não contaram com nenhum curso participante na região Nordeste.
As Faculdades foram representadas por dois cursos na região Nordeste, e os dois foram avaliados com conceito 3.
Na região Sudeste, as Universidades participaram com sete dos 22 cursos de Tecnologia em Agronegócio da região.
Entre os cursos em Universidades, nessa região, o conceito modal foi 4, atribuído a quatro cursos.
Os demais cursos receberam os conceitos 2 ( um curso ) e 3 ( dois cursos ).
Nenhum curso ficou sem conceito ( SC ) ou recebeu conceito 1 ou 5 Os Centros Universitários participaram com três cursos na região Sudeste, dos quais dois obtiveram o conceito modal 2, e um obteve conceito 3.
Já as Faculdades foram representadas por 12 cursos na região Sudeste, oito dos quais ficaram com conceito 3, conceito modal.
Aos demais cursos foram atribuídos o conceito 4 ( quatro cursos ).
Nenhum curso ficou sem conceito ( SC ) ou recebeu os demais conceitos restantes.
Dos 17 cursos da região Sul, nove eram de Universidades, para os quais o conceito modal foi 3, atribuído a quatro cursos.
Nesse tipo de organização, um curso ficou sem conceito ( SC ) e nenhum curso recebeu conceito 1 ou 2.
Os demais cursos receberam os conceitos 4 ( três cursos ) e 5 ( um curso ).
Os Centros Universitários participaram com um único curso na região Sul, este curso obteve conceito máximo, 5.
Quanto aos sete cursos vinculados a Faculdades na região Sul, quatro receberam os conceitos modais 4 e 5.
Os conceitos 1, 2 e 3 foram atribuídos a um curso cada.
Nenhum curso ficou sem conceito ( SC ).
Na região Centro-Oeste, seis dos 15 cursos eram de Universidades.
Nesse tipo de organização, nenhum curso recebeu conceito 1 ou ficou sem conceito ( SC ).
Os conceitos 2 e 3 foram atribuídos a dois cursos cada, conceitos modais.
Os demais conceitos, 4 e 5, foram atribuídos a um curso cada.
Os Centros Universitários não tiveram participação na região Centro-Oeste.
Dos nove cursos oferecidos em Faculdades na região Centro-Oeste, quatro receberam o

% matched lemmas: observação
\end{textsample}
